% Proved: closure of the same restricted alternating gate; four complete audits.
\subsection{Dimensional and finite-normalization closure of the alternating bound}

\begin{corollary}[Proved: alternating counts detected on fixed local surface groups]
\label{D:alternating-dimensional-normalization-closure}
Let $S$ be proper over an algebraically closed field of characteristic
$p>0$, and let $T\subset S$ be closed of codimension at least two
in every irreducible component that it meets. Assume either that
$T$ is finite, or that every normalized component of dimension at
least two has a projective modification satisfying the controlled
exceptional-centre hypothesis of
Theorem~\ref{D:generic-surface-controlled-boundary} for the inverse
image of $T$. Let $U$ be a connected component of $S\setminus T$,
and put $G=\pi_1(U)$.

There is a fixed finite collection of homomorphisms
\[
 \Lambda_{i,b}\longrightarrow G,\qquad
 \Lambda_{i,b}=\pi_1\bigl(\operatorname{Spec}R_{i,b}
                                      \setminus\{\mathfrak m_{i,b}\}\bigr),
                                                               \tag{D.ANC.1}
\]
where the $R_{i,b}$ are complete normal local surface rings with
algebraically closed coefficient fields. Their precise construction
is given below; they are fixed before $n$ or a cover is chosen.
For each group in (D.ANC.1), choose a finite collection of fixed
topologically finitely generated quotients. Call one of its orbits
on a finite $G$-set controlled if the orbit action factors through
one of these quotients. Fix $K>0$.

The isomorphism classes of degree-$n$ finite etale covers $X\to U$
such that
\[
 \sum_{i,b}\ 
 \sum_{\substack{O\in\Lambda_{i,b}\backslash X\\
                         O\text{ not controlled}}}|O|^2
                              \le K n\log(n+1)            \tag{D.ANC.2}
\]
are at most $n^{C n}$ for a constant $C$ depending only on the
fixed data. The covers in this assertion may be disconnected.
Consequently the distinct natural $A_n$ quotient kernels satisfying
(D.ANC.2) have multiplicity at most $|A_n|^c$ for one constant $c$,
and the corresponding selected alternating simple-product quotient
of $G$ is topologically finitely generated.

One may jointly choose the maximal pronilpotent and maximal
pro-prime-to-$p$ quotients. Thus (D.ANC.2) may charge only orbits
whose \emph{whole fixed local surface-group image} is neither
nilpotent nor of order prime to $p$, without a size restriction
on the other orbits. No condition is imposed merely on the smaller
curve-inertia orbits. Normality and reducedness of $S$ are not
required.
\end{corollary}

\begin{proof}
Reduce $S$ and normalize its irreducible components. Over $U$ this
gives a fixed finite proper surjection
\[
                         V=\coprod_i V_i\longrightarrow U
                                                               \tag{D.ANC.3}
\]
with connected normal $V_i$. For a component of dimension at least
two, choose once a normal proper surface $Z_i$ over an algebraically
closed extension of $k$, with a finite set $F_i$, and a surjection
\[
        \pi_1(Z_i\setminus F_i)\twoheadrightarrow\pi_1(V_i).
                                                               \tag{D.ANC.4}
\]
These are supplied by
Theorem~\ref{D:generic-surface-finite-punctures} in the finite case,
and by the codimension-at-least-two clause of
Theorem~\ref{D:generic-surface-controlled-boundary} in the other
case. In dimension two the original normal component itself may
be used. Apply
Lemma~\ref{D:finite-puncture-detecting-curve-local-types} to each
fixed pair $(Z_i,F_i)$. Its fixed completed-local groups are the
$\Lambda_{i,b}$ in (D.ANC.1); compose their maps with (D.ANC.4)
and the map from $V_i$ to $U$. Choose geometric base points and
connecting paths once. Changing those paths conjugates the maps
and changes neither orbit sizes nor factorization through the
chosen quotients. All the completed local rings here are attached
to these fixed \emph{detecting surfaces}; no assertion that they
are completions of the original higher-dimensional $S$ is made.

We first extend the surface count to disconnected covers. Keep one
fixed detecting curve $C^\circ=C\setminus B$, with $w=|B|>0$ and
$h=g(C)$, for a surface $(Z_i,F_i)$. The local-type enumeration
in the proof of
Corollary~\ref{D:fixed-local-quotient-alternating-count} counted
arbitrary finite local sets, so it already bounds the tuples of
full rank-$n$ local algebras under (D.ANC.2) by $n^{A_0 n}$ for
a fixed $A_0$. The discriminant sum is still at most $\lambda n$:
the upper-break bound applies to each actual cover and the
different estimate is additive over its local factors.

For a fixed tuple of full local algebras, consider a cover of
$C^\circ$ with $k\le n$ connected components. Temporarily order
those components. At a marked point its completed algebra has
at most $n$ connected factors. Assigning those factors to the
$k$ labelled global components has at most $n^n$ possibilities.
Doing this at every marked point, and allowing the $n$ choices of
$k$, costs at most $n^{w n+1}$. Assignments with inconsistent
degrees or empty global components are discarded. Since $B$ is
nonempty, every actual connected component has positive degree
at each marked point and is included in this enumeration.

For a consistent assignment, let $d_j$ and $D_j$ be the degrees
and discriminant sums of its connected components. Then
$\sum_jd_j=n$ and $\sum_jD_j\le\lambda n$. The connected
fixed-formal-fibre bound in
Theorem~\ref{D:fixed-formal-curve-fibre-bound} gives at most
\[
 \prod_j(d_j!)^{2h}d_j^{2D_j}
                  \le (n!)^{2h}n^{2\lambda n}             \tag{D.ANC.5}
\]
such ordered tuples of global covers. Forgetting the ordering
can only lower the count. This proves a bound $n^{A_i n}$ for
all degree-$n$ covers of $Z_i\setminus F_i$ satisfying the local
condition, with $A_i$ fixed. The original detecting-curve
surjection is injective on finite-set isomorphism classes even
for nontransitive actions: an equivariant bijection after
precomposition is equivariant for the entire quotient group.
The same applies to (D.ANC.4), giving the bound for the relevant
degree-$n$ covers of $V_i$.

Components of dimension at most one meet no omitted point, by
the codimension hypothesis. They are proper normal curves or
points. Their fundamental groups have fixed finite generating
sets by \citet[Expose~X, Theorem~2.6, pp.~272--273]{C_SGA1}.
If $d_i$ elements suffice, all their degree-$n$ covers have at
most $(n!)^{d_i}\le n^{d_i n}$ isomorphism classes. They need
no groups (D.ANC.1) and no energy charge.

Fix one tuple of rank-$n$ covers on the $V_i$. On
$V\times_U V$, the finite etale Isom scheme between the two
pullbacks has degree $n!$. A section on a connected component
is determined by its value at one geometric point. Thus there
are at most $(n!)^a$ descent isomorphisms, where the fixed integer
$a=\#\pi_0(V\times_U V)$ is independent of $n$. The cocycle
condition and passage to isomorphism classes only decrease the
count. Finite proper descent of finite etale covers is effective
and fully faithful
\citep[Expose~IX, Theorem~4.7]{C_SGA1}. Multiplying the fixed
component bounds by $(n!)^a\le n^{a n}$ proves the asserted
$n^{Cn}$ bound. Rank one has only the trivial cover of connected
$U$, so the stated form also holds for $n=1$.

The selected alternating-kernel and simple-product conclusions
now follow exactly as in
Theorem~\ref{D:finite-puncture-local-orbit-factorial-count}:
distinct kernels give distinct natural permutation classes,
$|A_n|\ge n^{n/8}$ for $n\ge5$, and the uniform polynomial
multiplicity bound gives one generating bound by the cited
LMT estimate, Goursat's lemma and compactness.

For the concrete final choice, each ring in (D.ANC.1) has a
finitely generated maximal pronilpotent quotient by
Corollary~\ref{D:fixed-local-quotient-alternating-count}. The other
quotient is also a \emph{known} finitely generated group:
\citet[\S9.1 and Theorem~9.2, pp.~122--123]
{C_GrothendieckMurre1971PrimeToP} prove finite presentation of the
maximal prime-to-$p$ quotient for a connected open in the spectrum
of a complete noetherian local ring of dimension two with
algebraically closed residue field. The finite presentation is in
the pro-prime-to-$p$ category; its stated definition includes a
finite set of topological generators, which is all that is used
here. Every ring in (D.ANC.1) satisfies those hypotheses. A finite
action factors through this quotient if and only if its image
has order prime to $p$. Choosing these two fixed quotients
therefore controls precisely the union of the nilpotent-image
and prime-to-$p$-image orbits, as asserted.
\end{proof}

\begin{remark}
This is a closure of the same restricted alternating count.
The local homomorphisms and quotients are part of the fixed
detecting data. The conclusion does not discard unrestricted
nonsolvable local images, and does not prove an unrestricted
factorial cover bound or finite generation of the whole $G$.
\end{remark}
